% TOP_PROBLEM: {"id": 1469, "title": "Serre's Positivity Conjecture", "category_id": 4, "category": {"id": 4, "name": "algebra", "display_name": "Algebra", "description": "Group theory, ring theory, field theory, and algebraic structures.", "slug": "algebra", "order_index": 4, "created_at": "2026-07-31T15:26:25.671Z"}, "status": "open"}
% FIRST_POSTED: 2026-10-06
% ATTEMPT_STATUS: solved
% OPENAI_CLAIM: {"schema_version":1,"source_repo":"https://github.com/openai/math","source_commit":"adc7f1241b42e322a6451854ab7e4b4c146bf78a","release_date":"2026-10-06","match":"full","resolution":"proved","independently_reviewed":false,"second_pass":"agreed","families":[{"family":"193","title":"Serre’s intersection-multiplicity conjecture","subject":"Algebra","manuscripts":[{"title":"Positivity of Serre's Intersection Multiplicity","pdf_path":"preprints/Positivity-of-Serres-Intersection-Multiplicity-September-23-2026/paper.pdf","readme_path":"preprints/Positivity-of-Serres-Intersection-Multiplicity-September-23-2026/README.md","date":"2026-09-23","theorem_ref":"Theorem 1.1"}],"lean":null,"reasoning_trace":null}]}
\documentclass[11pt]{article}
\usepackage[margin=25mm]{geometry}
\usepackage{fontspec}
\setmainfont{Latin Modern Roman}
\usepackage{amsmath,amssymb,mathtools}
\usepackage{unicode-math}
\setmathfont{Latin Modern Math}
\usepackage{enumitem}
\usepackage{xurl,hyperref,bookmark}
\hypersetup{colorlinks=true,urlcolor=blue}
\setcounter{secnumdepth}{0}
\setlength{\parindent}{0pt}
\setlength{\parskip}{5pt}
\setlength{\emergencystretch}{3em}
\begin{document}
\section*{Serre's Positivity Conjecture}

\subsection{OpenAI's claim}
OpenAI presents a proof in the selected manuscript. Serre intersection numbers are claimed to be positive when two nonzero finite modules have complementary support dimensions and a finite-length tensor product over a regular local ring, even in ramified mixed characteristic.

\subsection{Outline of the argument}
The reductions pass to completed quotient domains and algebras carrying normalized length. Galois-averaged supported K-theory compares ordinary Tor Euler numbers with normalized ones; parameter Koszul exactness and truncated finite resolutions then force each relevant normalized contribution to be positive.

\subsection{Scope relative to this problem}
Theorem 1.1 has exactly the record’s regular-local, finite-length and complementary-dimension hypotheses and proves strict positivity without characteristic or ramification restrictions. Within these assumptions, this is a claimed affirmative answer to the stated question.

\subsection{Formal verification}
Family 193 has no Lean scope page in this release. Its selected result is not presented as formalized. This site has not executed the Lean code. The scope descriptions do not certify the remaining manuscript arguments.

\subsection{Status caveat}
OpenAI attributes these manuscripts to work by an internal model. Its release cautions that arguments without formalization can contain mistakes. This catalogue has not independently checked the proofs; the labels report the mathematical claims attributed to OpenAI.

\subsection{Sources}
\begin{itemize}
\item OpenAI, \emph{Positivity of Serre's Intersection Multiplicity}, 2026-09-23; citation key \texttt{OAI:Positivity-of-Serres-Intersection-Multiplicity-September-23-2026}. \href{https://github.com/openai/math/blob/adc7f1241b42e322a6451854ab7e4b4c146bf78a/preprints/Positivity-of-Serres-Intersection-Multiplicity-September-23-2026/paper.pdf}{PDF (GitHub)}; \href{https://github.com/openai/math/raw/adc7f1241b42e322a6451854ab7e4b4c146bf78a/preprints/Positivity-of-Serres-Intersection-Multiplicity-September-23-2026/paper.pdf}{PDF (direct)}; \href{https://github.com/openai/math/blob/adc7f1241b42e322a6451854ab7e4b4c146bf78a/preprints/Positivity-of-Serres-Intersection-Multiplicity-September-23-2026/README.md}{Manuscript page and citation}; \href{https://github.com/openai/math/blob/main/preprints/Positivity-of-Serres-Intersection-Multiplicity-September-23-2026/paper.pdf}{Latest version}.
\item \href{https://github.com/openai/math/blob/adc7f1241b42e322a6451854ab7e4b4c146bf78a/CONTENTS.md}{OpenAI catalogue entry}.
\end{itemize}
\end{document}
